% Part: first-order-logic
% Chapter: axiomatic-deduction
% Section: deduction-theorem

% verification of properties of provability needed for maximally
% consistent sets in the completeness chapter.

\documentclass[../../../include/open-logic-section]{subfiles}

\begin{document}

\iftag{FOL}
      {\olfileid{fol}{axd}{ded}}
      {\olfileid{pl}{axd}{ded}}

\olsection{निगमन प्रमेय}

आपण पाहिल्याप्रमाणे, स्वयंसिद्धकीय पद्धतीत !!{derivation}s देणे किचकट
असते आणि !!{derivation}s शोधणे कठीण असू शकते. !!{formula}s च्या लांब
याद्या प्रत्यक्ष लिहिण्याऐवजी काही सोपे निष्कर्ष वापरून अशा
!!{derivation}s अस्तित्वात आहेत असा युक्तिवाद करणे सहसा सोपे असते. असे तीन
निष्कर्ष आपण आधीच प्रस्थापित केले आहेत: $!A \in \Gamma$ हे माहीत असेल,
तर $\Gamma \Proves !A$ असे नेहमी म्हणता येते, हे
\olref[ptn]{prop:reflexivity} सांगतो. $\Gamma \Proves !A$ असेल, तर
$\Gamma \cup \{!B\} \Proves !A$ सुद्धा असते, हे
\olref[ptn]{prop:monotonicity} सांगतो. आणि $\Gamma \Proves !A$ व
$!A \Proves !B$ असेल, तर $\Gamma \Proves !B$, असे
\olref[ptn]{prop:transitivity} वरून मिळते. पुढे आणखी एक सोपा निष्कर्ष,
विध्यात्मक अनुमानाची ``अधि''-आवृत्ती, दिली आहे:

\begin{prop}
  \ollabel{prop:mp} जर $\Gamma \Proves !A$ आणि $\Gamma \Proves !A \lif
  !B$, तर $\Gamma \Proves !B$.
\end{prop}

\begin{proof}
  $\{!A, !A \lif !B\} \Proves !B$ हे आपल्याकडे आहे:
  \begin{derivation}
    1. & $!A$ & Hyp.\\
    2. & $!A \lif !B$ & Hyp.\\
    3. & $!B$ & 1, 2, MP
  \end{derivation}
  \olref[ptn]{prop:transitivity} वरून $\Gamma \Proves !B$.
\end{proof}

या संदर्भात आपण वापरणार असलेला सर्वांत महत्त्वाचा निष्कर्ष म्हणजे निगमन
प्रमेय:

\begin{thm}[निगमन प्रमेय]
\ollabel{thm:deduction-thm} $\Gamma \cup \{!A\} \Proves !B$ तेव्हा आणि
केवळ तेव्हाच, जेव्हा $\Gamma \Proves !A \lif !B$.
\end{thm}

\begin{proof}
``जर'' ही दिशा तात्काळ मिळते. $\Gamma \Proves !A \lif !B$ असेल, तर
\olref[ptn]{prop:monotonicity} वरून $\Gamma \cup \{!A\} \Proves !A \lif !B$
सुद्धा असते. तसेच \olref[ptn]{prop:reflexivity} वरून
$\Gamma \cup \{!A\} \Proves !A$. म्हणून \olref{prop:mp} वरून
$\Gamma \cup \{!A\} \Proves !B$.

``केवळ जर'' या दिशेसाठी, $\Gamma \cup \{!A\}$ पासून $!B$ च्या
!!{derivation} च्या लांबीवर विगमन करू.

विगमनाच्या आधारपायरीसाठी, लांबी~$1$ असलेल्या प्रत्येक !!{derivation} साठी
दावा सिद्ध करू. $\Gamma \cup \{!A\}$ पासून~$!B$ ची लांबी~$1$ असलेली
!!^a{derivation} म्हणजे एकटे $!B$; ती योग्य असेल, तर $!B$ हे एकतर
$\in \Gamma \cup \{!A\}$ असते किंवा स्वयंसिद्धक असते. $!B \in \Gamma$
असेल किंवा ते स्वयंसिद्धक असेल, तर $\Gamma \Proves !B$. तसेच
\olref[prp]{ax:lif1} वरून $\Gamma \Proves !B \lif (!A \lif !B)$, आणि
\olref{prop:mp} वरून $\Gamma \Proves !A \lif !B$. $!B \in \{ !A\}$
असेल, तर $\Gamma \Proves !A \lif !B$, कारण शेवटचे !!{sentence}~$!A
\lif !B$ हे $!A \lif !A$ सारखेच आहे आणि ते आपण
\olref[pro]{ex:identity} मध्ये !!{derive}d केले आहे.

विगमन पायरीसाठी, $\Gamma \cup \{!A\}$ पासून~$!B$ ची एखादी
!!a{derivation} विध्यात्मक अनुमानाने समर्थित असलेल्या~$!B$ या पायरीने
संपते असे समजा. (तिला विध्यात्मक अनुमानाने समर्थन मिळत नसेल, तर
$!B \in \Gamma$, $!B \ident !A$, किंवा $!B$ हे स्वयंसिद्धक असते, आणि
विगमनाच्या आधारपायरीतील युक्तिवादच लागू होतो.) मग एखाद्या
!!{formula}~$!C$ साठी त्या !!{derivation} मधील आधीच्या काही पायऱ्या
$!C \lif !B$ आणि $!C$ असतात; म्हणजे $\Gamma \cup \{!A\} \Proves !C \lif !B$
आणि $\Gamma \cup \{!A\} \Proves !C$. त्यांच्याशी संबंधित
!!{derivation}s लहान असल्यामुळे त्यांना विगमन गृहीतक लागू होते. म्हणून
आपल्याकडे पुढील दोन्ही आहेत:
\begin{align*}
  & \Gamma \Proves !A \lif (!C \lif !B); \\
  & \Gamma \Proves !A \lif !C.
\end{align*}
पण \olref[prp]{ax:lif2} वरून पुढीलही मिळते:
\[
\Gamma \Proves (!A \lif (!C \lif !B)) \lif
((!A\lif !C)  \lif (!A \lif !B)),
\]
आणि \olref{prop:mp} चे दोन उपयोग केल्यावर आवश्यकतेनुसार
$\Gamma \Proves !A \lif !B$ मिळते.
\end{proof}

निगमन प्रमेय खरे ठरावे यासाठीच \olref[prp]{ax:lif1} आणि
\olref[prp]{ax:lif2} यांची नेमकी अशी निवड केली होती, हे लक्षात घ्या.

!!{derivability} विषयीची पुढील काही उपयुक्त तथ्ये सरावासाठी सोडली आहेत.

\begin{prop}
  \ollabel{prop:derivfacts}
\begin{enumerate}
\item $\Proves (!A \lif !B) \lif ((!B \lif !C)
  \lif (!A \lif !C))$; \ollabel{derivfacts:a}
%READERNOTE{OLAXD-001: गोठवलेल्या इंग्रजी सूत्रात सर्वांत बाहेरचा उजवा कंस सुटला आहे. मराठी सूत्रात तो कंस जोडला असून गोठवलेले इंग्रजी बाइट्स बदललेले नाहीत.}
\item $\Gamma \cup \{ \lnot !A\}
  \Proves \lnot !B$ असेल, तर $\Gamma \cup \{ !B\} \Proves
  !A$ (प्रतिपरिवर्तन); \ollabel{derivfacts:b}
\item  $\{ !A, \lnot!A\} \Proves
    !B$ (Ex Falso Quodlibet, म्हणजे असत्यापासून काहीही; स्फोट); \ollabel{derivfacts:c}
\item  $\{ \lnot\lnot!A\} \Proves
  !A$ (द्विनकरण विलोपन);\ollabel{derivfacts:d}
\item $\Gamma \Proves \lnot\lnot!A$ असेल, तर $\Gamma \Proves
  !A$;\ollabel{derivfacts:e}
\end{enumerate}
\end{prop}

\tagprob{FOL}
\begin{prob}
  \olref[fol][axd][ded]{prop:derivfacts} सिद्ध करा.
\end{prob}
\tagendprob

\tagprob{notFOL}
\begin{prob}
  \olref[pl][axd][ded]{prop:derivfacts} सिद्ध करा.
\end{prob}
\tagendprob

\end{document}
